\documentclass[../../../include/open-logic-section]{subfiles}

\begin{document}
\olfileid{sth}{card-arithmetic}{simp}

\olsection{د جمعې او ضرب ساده کول}

د اصلي عددونو نامتناهي جمع او ضرب په پايله کښې \emph{ډېر} اسانه
ثابتېږي۔ دا له دې راځي چې اصلي عددونه (ځينې) ترتيبي عددونه دي،
نو ښه ترتيب شوي دي، او له همدې امله په يوه ټاکلې توګه کار پرې
کېدے شي۔ خو د دې ښودل \emph{هغومره} اسانه نه دي۔ د پيل دپاره
يو څه ځيرک تعريف ته اړتيا لرو:

\begin{defn}
پر ترتيبي عددونو جوړو يو \emph{معياري ترتيب}، $\canonord$،
داسې تعريفوو چې $\tuple{\alpha_1, \alpha_2} \canonord
\tuple{\beta_1, \beta_2}$ هغه وخت او يوازې هغه وخت وي چې له
لاندې حالتونو يو پېښ شي:
\begin{enumerate}
	\item $\max(\alpha_1, \alpha_2) < \max(\beta_1, \beta_2)$؛ يا
	\item $\max(\alpha_1, \alpha_2) = \max(\beta_1, \beta_2)$ او
	$\alpha_1 < \beta_1$؛ يا
	\item $\max(\alpha_1, \alpha_2) = \max(\beta_1, \beta_2)$ او
	$\alpha_1 = \beta_1$ او $\alpha_2 < \beta_2$
\end{enumerate}
\end{defn}

\begin{lem}
$\tuple{\alpha \times \alpha, \canonord}$ د هر ترتيبي عدد
$\alpha$ دپاره ښه ترتيب دے۔
\end{lem}

\begin{proof}
څرګنده ده چې $\canonord$ پر $\alpha \times \alpha$ هرې دوو
بېلابېلو جوړو ته ترتيب ورکوي۔ فرض کړئ نه $\tuple{\alpha_1,
\alpha_2}$ له $\tuple{\beta_1,
\beta_2}$ څخه د $\canonord$ له مخې کوچنۍ وي، نه دويمه له
لومړۍ څخه۔ بيا $\max(\alpha_1,
\alpha_2) = \max(\beta_1, \beta_2)$، $\alpha_1 = \beta_1$
او $\alpha_2 = \beta_2$؛ نو $\tuple{\alpha_1, \alpha_2} =
\tuple{\beta_1,  \beta_2}$۔

د ښه ترتيب د ښودلو دپاره دې $X \subseteq \alpha\times\alpha$
ناتش وي۔ څرنګه چې $\alpha$ ترتيبي عدد دے، يو $\delta$ د
$\Setabs{\max(\gamma_1, \gamma_2)}{\tuple{\gamma_1,
\gamma_2} \in X}$ تر ټولو کوچنے غړے دے۔ اوس له
$\Setabs{\tuple{\gamma_1,\gamma_2} \in X}{\max(\gamma_1, \gamma_2) =
\delta}$ څخه ټولې جوړې لرې کړئ، پرته له هغو چې لومړے مختص يې
تر ټولو کوچنے وي؛ له پاتې جوړو څخه هغه چې دويم مختص يې تر
ټولو کوچنے وي، د $\canonord$ له مخې د $X$ تر ټولو کوچنے غړے دے۔
\end{proof}
\noindent
اوس يوې وړوکې، ساده کتنې ته راځو:

\begin{prop}\ollabel{simplecardproduct}
که $\cardeq{\alpha}{\beta}$ وي، نو $\cardeq{\alpha \times \alpha}{\beta
\times \beta}$۔
\end{prop}

\begin{proof}
بس $f \colon \alpha \to \beta$ دې دا تابع وټاکي:
$\tuple{\gamma_1,
\gamma_2} \mapsto \tuple{f(\gamma_1), f(\gamma_2)}$۔
\end{proof}
\noindent
اوس دا ټول د يوه مهمې لمې په ثبوت کښې کاروو:
\begin{lem}\ollabel{alphatimesalpha}
$\cardeq{\alpha}{\alpha \times \alpha}$، د هر نامتناهي ترتيبي
عدد $\alpha$ دپاره۔
\end{lem}

\begin{proof}
د تناقض دپاره دې $\alpha$ تر ټولو کوچنے نامتناهي ترتيبي عدد
وي چې دا ادعا يې په اړه ناسمه وي۔
\olref[sfr][siz][zigzag]{natsquaredenumerable} ښيي چې
$\cardeq{\omega}{\omega\times\omega}$، نو $\omega \in \alpha$۔
سربېره پر دې، $\alpha$ اصلي عدد دے: د تناقض دپاره فرض کړئ
چې نه دے؛ بيا $\card{\alpha} \in \alpha$، نو د فرضيې له مخې
$\cardeq{\card{\alpha}}{\card{\alpha} \times \card{\alpha}}$؛
او د تعريف له مخې $\cardeq{\card{\alpha}}{\alpha}$؛ نو له
\olref{simplecardproduct} څخه
$\cardeq{\alpha}{\alpha\times\alpha}$ راځي۔

اوس د هر $\tuple{\gamma_1, \gamma_2} \in \alpha \times \alpha$
دپاره لاندې برخه وګورئ:
\begin{align*}
	\text{Seg}(\gamma_1, \gamma_2) &= \Setabs{\tuple{\delta_1, \delta_2} \in \alpha \times \alpha}{\tuple{\delta_1, \delta_2} \canonord \tuple{\gamma_1, \gamma_2}}
\end{align*}
که $\gamma = \max(\gamma_1, \gamma_2)$ ونيسو، نو
$\tuple{\gamma_1, \gamma_2} \canonord \tuple{\gamma+1, \gamma + 1}$۔
نو کله چې $\gamma$ نامتناهي وي، لاندې ته پام وکړئ:
\begin{align*}
	\text{Seg}(\gamma_1, \gamma_2) &
	\precsim ((\gamma \ordplus 1)\ordtimes (\gamma \ordplus 1))\\
	&\approx (\gamma \ordtimes \gamma)
	\text{، د \olref[ord-arithmetic][using-addition]{ordinfinitycharacter} او
	\olref{simplecardproduct} له مخې}\\
	&\approx \gamma \text{، د استقرا د فرضيې له مخې}\\
	& \prec \alpha\text{، ځکه $\alpha$ اصلي عدد دے}
\end{align*}
نو $\ordtype{\alpha\times \alpha, \canonord} \leq \alpha$،
او له دې $\cardle{\alpha \times \alpha}{\alpha}$۔ څرنګه چې
طبعاً $\cardle{\alpha}{\alpha \times \alpha}$، پايله د
شرودر--برنشتاين له قضيې راځي۔
\end{proof}

بالاخره د ساده کولو پايلې ته رسېږو:

\begin{thm}\ollabel{cardplustimesmax}
که $\cardfont{a}, \cardfont{b}$ نامتناهي اصلي عددونه وي، نو:
\[
	\cardfont{a}
\cardtimes \cardfont{b} = \cardfont{a} \cardplus \cardfont{b} =
\text{max}(\cardfont{a}, \cardfont{b}).
\]
\end{thm}

\begin{proof}
له عموميت له کمېدو پرته فرض کړئ $\cardfont{a} = \max(\cardfont{a},
\cardfont{b})$۔ بيا د \olref{alphatimesalpha} په کارولو سره
$\cardfont{a}\cardtimes\cardfont{a} = \cardfont{a} \leq \cardfont{a}
\cardplus \cardfont{b} \leq \cardfont{a} \cardplus \cardfont{a} \leq
\cardfont{a} \cardtimes \cardfont{a}$۔ \end{proof}\noindent په ورته ډول،
که $\cardfont{a}$ نامتناهي وي، د $\cardfont{a}$ په اندازه د
داسې سټونو اتحاد چې د هر يوه اندازه $\leq\cardfont{a}$ وي،
اندازه يې هم $\leq\cardfont{a}$ وي:

\begin{prop}\ollabel{kappaunionkappasize}
دې $\cardfont{a}$ يو نامتناهي اصلي عدد وي۔ د هر ترتيبي عدد
$\beta \in \cardfont{a}$ دپاره دې $X_\beta$ داسې سټ وي چې
$\card{X_\beta} \leq \cardfont{a}$۔ بيا
$\card{\bigcup_{\beta \in \cardfont{a}} X_\beta}
\leq \cardfont{a}$۔
\end{prop}

\begin{proof}
د هر $\beta \in \cardfont{a}$ دپاره يوه !!a{injection}
$f_\beta \colon X_\beta \to \cardfont{a}$ وټاکئ۔\footnote{دا تابعې
څنګه يوځاے «ټاکل» کېږي؟ \olref[sth][choice][countablechoice]{sec}
وګورئ۔} يوه !!a{injection} $g \colon
\bigcup_{\beta \in \cardfont{a}} X_\beta \to \cardfont{a} \times
\cardfont{a}$ په $g(v) = \tuple{\beta, f_\beta(v)}$ سره تعريف
کړئ، چې دلته $v \in X_\beta$، او $v \notin X_\gamma$ د هر
$\gamma \in \beta$ دپاره۔ اوس
$\bigcup_{\beta \in \cardfont{a}} X_\beta \preceq \cardfont{a} \times
\cardfont{a} \approx \cardfont{a}$ د \olref{cardplustimesmax}
له مخې۔
\end{proof}

\end{document}
